% ECONOMICS_PROBLEM: {"id": "C21-173", "title": "Identification and inference after learning a mediator representation", "jel_code": "C21", "source_id": "OP-0517"}
% !TeX program = xelatex
\documentclass[11pt,a4paper]{article}
\usepackage[margin=23mm]{geometry}
\usepackage{fontspec}
\setmainfont{Latin Modern Roman}
\usepackage{amsmath,amssymb,mathtools}
\usepackage{xcolor,enumitem,booktabs,longtable,array,xurl,hyperref}
\definecolor{muted}{HTML}{53636C}
\hypersetup{colorlinks=true,urlcolor=blue,linkcolor=blue}
\setlength{\parindent}{0pt}
\setlength{\parskip}{5pt}
\newcommand{\R}{\mathbb R}
\newcommand{\E}{\mathbb E}
\newcommand{\N}{\mathbb N}
\newcommand{\ind}{\mathbf 1}
\newcommand{\Var}{\operatorname{Var}}
\newcommand{\Cov}{\operatorname{Cov}}
\newcommand{\supp}{\operatorname{supp}}
\DeclareMathOperator*{\argmin}{arg\,min}
\DeclareMathOperator*{\argmax}{arg\,max}
\newcommand{\entrymeta}[3]{{\sffamily\small\color{muted}\textbf{\texttt{#1}}\quad|\quad #2\par #3\par}}
\newcommand{\Problemref}[1]{\texttt{#1}}
\newcommand{\JELref}[2]{\texttt{#2} (JEL \texttt{#1})}
\begin{document}
\section*{Identification and inference after learning a mediator representation}
\textbf{Problem ID:} \texttt{C21-173}\quad \textbf{JEL:} \texttt{C21}\par
% BEGIN ECONOMICS STATEMENT
\label{problem:OP-0517}
\label{beyond:CI-38}
\entrymeta{CI-38}{Econometric theory and causal inference}{Research agenda formalization}
\subsection*{Definitions and assumptions}
Specify jointly measurable potential outcomes $Y(w,m)=F(w,m,X,U)$ with $U\perp(W,M)\mid X$, consistency, no interference, and intervention-support positivity. Let $X\in[0,1]^d$, $M\in[0,1]^p$, and let $\Phi$ have uniform-metric covering numbers satisfying $\log N(\epsilon,\Phi)\leq v\log(A/\epsilon)$ for $0<\epsilon\leq1$. Assume some $\phi_*\in\Phi$ yields the smooth reduced causal response. Independent training selects $\widehat\phi$ with $\Pr(\|\widehat\phi-\phi_*\|_\infty>a_n)\leq\delta_n$ for specified sequences $a_n,\delta_n\to0$. The induced full-mediator intervention must be absolutely continuous with a bounded density ratio relative to the observed mediator law. A representation intervention requires a lift: for every $\phi$, specify a feasible kernel $K_\phi(dm\mid z,x)$ supported on $\{m:\phi(m)=z\}$. With representation draw $g(dz\mid x)$ and lift randomization independent of $U$ conditional on $X$, the controlled mean is
\[
 \Psi_{\phi,K}(w,g)=
 \E_X\int\!\int\mu_w(x,m)K_\phi(dm\mid z,X)g(dz\mid X),
\]
where $\mu_w(x,m)=\E[Y(w,m)\mid X=x]$.
\subsection*{Formal question}
Characterize when $\Psi_{\phi,K}$ is invariant to every allowed lift $K_\phi$, using explicit restrictions on within-fiber causal responses and feasible supports. For approximate invariance, bound the identified range when within-fiber oscillations of $\mu_w$ are at most $\delta$. Then construct confidence regions for $\Psi_{\widehat\phi,K}$ conditional on the training sample, using independent inference data. Determine attainable rates under reduced responses in $\mathcal H_{d+r}^\beta(L)$, reduced mediator densities in $\mathcal H_{d+r}^\gamma(L)$, propensity in $\mathcal H_d^\alpha(L)$, representation complexity, and a specified training-error bound; distinguish uncertainty about the learned target from uncertainty about a fixed scientific target.
\subsection*{Scope and status}
This lift-based question is an editorial refinement of the representation-learning agenda. Merely conditioning on an embedding does not define an intervention on it. A representation selected to predict the observed outcome need not yield invariant causal meaning, and sample splitting addresses statistical selection rather than this conceptual defect. If lifts induce different outcomes, the effect must be indexed by a lift or reported as a set instead of being called uniquely identified.
\subsection*{References for the source of the problem}
Carlos Cinelli, Avi Feller, Guido Imbens, Edward Kennedy, Sara Magliacane, and Jose Zubizarreta, \emph{Challenges in Statistics: A Dozen Challenges in Causality and Causal Inference}, arXiv:2508.17099v1 [stat.ME], 23 August 2025, \S3.4.3, ``Representing complex mediators,'' discussion of representation learning. \url{https://arxiv.org/html/2508.17099v1#S3.SS4.SSS3}.

\subsection*{Common conventions: Source mathematical conventions}

\textbf{Well-defined objects.}
All probability spaces and measurable variables are specified or inherited
from a local statistical model. Expectations used as finite targets require
the stated integrability. A decision rule or estimator is measurable with
respect to its available information; randomized rules may use an auxiliary
independent random variable. Norms without a subscript are Euclidean in
finite-dimensional spaces. An optimizer is requested only when attainment
is assured; otherwise the question concerns the optimal value and
$\varepsilon$-optimal admissible rules for every $\varepsilon>0$.

\textbf{Identification.}
A structural model $m\in\mathcal M$ implies an observable law
$P_m^{\rm obs}$ and target $\theta(m)$. The target is point identified on
$\mathcal M$ if
\[
 P_{m_1}^{\rm obs}=P_{m_2}^{\rm obs}
 \quad\Longrightarrow\quad\theta(m_1)=\theta(m_2)
 \qquad(m_1,m_2\in\mathcal M).
\]
When this fails, the sharp identified set is
\[
 \Theta(P;\mathcal M)=\{\theta(m):m\in\mathcal M,
                                  P_m^{\rm obs}=P\}.
\]
Specifying an intervention or estimand does not establish identification.
Positivity, exclusion, causal sufficiency, measurement restrictions, and
transport assumptions are substantive local inputs, never automatic
consequences of notation.

\textbf{Minimax risk and nontrivial inference.}
For a statistical experiment with data law $P^{(n)}$, target $\theta(P)$,
and nonnegative loss $\ell$, define
\[
 \mathcal R_n(\mathcal P,\ell)
   =\inf_{\widehat\theta_n}\sup_{P\in\mathcal P}
      \E_{P^{(n)}}\ell(\widehat\theta_n,\theta(P)).
\]
The infimum is over measurable estimators. A sharp rate is a sequence
$r_n$ for which $0<c\leq C<\infty$, independent of $n$, satisfies
$c r_n\leq\mathcal R_n\leq C r_n$ for all sufficiently large $n$.
Squared-error parametric risk is of order $n^{-1}$; the corresponding
root-mean-squared error is of order $n^{-1/2}$. These different conventions
are distinguished locally. A uniform asymptotic confidence region $C_n$
must satisfy
\[
 \liminf_{n\to\infty}\inf_{P\in\mathcal P}
     \Pr_{P^{(n)}}\{\theta(P)\in C_n\}\geq 1-\alpha,
 \qquad 0<\alpha<1,
\]
together with an informative diameter or expected-width requirement.
Coverage alone permits the vacuous whole parameter space. Testing questions
similarly impose both size control and power against a declared separated
alternative class.

\textbf{Binary causal baseline.}
When an entry explicitly invokes this baseline, one observes independent
copies of $O=(X,W,Y)$, with $W\in\{0,1\}$ and potential outcomes $Y(0),Y(1)$.
Assume consistency $Y=Y(W)$, no interference, conditional exchangeability
$(Y(0),Y(1))\perp W\mid X$, and overlap
$\epsilon\leq e(X)\leq1-\epsilon$ almost surely for a fixed
$0<\epsilon<1/2$, where
\[
 e(x)=\Pr(W=1\mid X=x),\qquad
 m_w(x)=\E[Y\mid W=w,X=x],\qquad
 \tau(x)=m_1(x)-m_0(x).
\]
These assumptions identify conditional mean effects almost everywhere
under the covariate law. Pointwise or uniform effects require appropriate
regularity and a specified support. Continuous treatment, longitudinal
data, adaptive experiments, and network interference require their own
assumptions in place of this baseline. Bounded outcomes, densities, and
smoothness are additional assumptions only where stated.

\textbf{Function classes.}
On $[0,1]^d$, $\mathcal H_d^s(L)$ is the isotropic Hölder ball of radius
$L>0$: put $k=\lceil s\rceil-1$ and $\gamma=s-k\in(0,1]$, bound derivatives
through order $k$, and require the order-$k$ derivatives to be
$\gamma$-Hölder with constant at most $L$. Thus integer orders use a
Lipschitz final derivative convention. The class
$\operatorname{BV}_d(L)$ consists of functions $f\in L^1((0,1)^d)$ whose
distributional derivative is a finite vector measure and for which
$\|f\|_{L^1}+|Df|((0,1)^d)\leq L$.
An $L^\infty$ bound or a particular pointwise version is imposed separately
when needed. Sparse and neural-network classes require a declared
dictionary or architecture and coefficient bounds.

An anisotropic H\"older class with declared block smoothness indices means
coordinatewise H\"older regularity in each block, uniformly in the other
blocks, with all corresponding derivative and H\"older bounds at most
the stated radius. Derivative orders and the integer-order convention in
each block are those just given. Mixed-derivative bounds are additional
restrictions only when explicitly stated.

\textbf{Decision and computational questions.}
Policy regret compares admissible feasible policies under the same law and
constraints. In computational entries, finite data and thresholds have
rational encodings; running time is measured in their total bit length.
Real-valued equilibrium search uses the explicitly declared algebraic or
FIXP encoding. An approximation ratio for maximization is written
$\operatorname{OPT}/\operatorname{ALG}$ and for minimization as
$\operatorname{ALG}/\operatorname{OPT}$, with zero optima and infeasible
instances handled locally. Historical approximation bounds are not
automatically current bounds.
% END ECONOMICS STATEMENT
\end{document}
